\section*{第 \ref{ch:g7:light-sources-propagation} 章——光源与光的传播}

\begin{solution}{exo:g7:light-sources-propagation:1}
\omterm{def:g7:light-sources-propagation:primary}{自身发光的光源}：烛焰、萤火虫、红色交通
信号灯。\omterm{def:g7:light-sources-propagation:diffusion}{漫射物体}：\omterm{def:g2:moon-in-the-sky:moon}{月球}、这一页、金星——
\omterm{def:g4:solar-system:planet}{行星}靠漫射的阳光发亮，无论它看上去多像\omterm{def:g4:solar-system:star}{恒星}。
\end{solution}

\begin{solution}{exo:g7:light-sources-propagation:2}
它将接收到的部分光向各个方向重新发出。在
暗室中，它接收不到任何光，也转发不出任何光——
因而看不见。
\end{solution}

\begin{solution}{exo:g7:light-sources-propagation:3}
光必须从物体射向眼睛。完美的抛光面
拒绝\emph{漫射}：它只把光沿一个规整的方向送出，
因而没有光向每个视角宣告“这里是表面”，
眼睛接收到的只是其他物体经它改道而来的景象，
即其他物体的像。
\end{solution}

\begin{solution}{exo:g7:light-sources-propagation:4}
苹果同时向各个方向漫射——每只眼睛都
能分到一份。勺子的抛光面把灯光沿一个方向送出，
因此只有一个眼睛位置处在这条路上。
\end{solution}

\begin{solution}{exo:g7:light-sources-propagation:5}
洁净的\omterm{def:g2:air-around-us:air}{空气}几乎不向侧面转发光，因此\omterm{def:g6:light-propagation:ray}{光束}中的光
不会射入\omterm{def:g6:light-propagation:ray}{光束}路径之外的眼睛。薄雾中布满细小的
\omterm{def:g7:light-sources-propagation:diffusion}{漫射物体}，每个都把一部分光散射向你：是漫射
划出了这条界线。
\end{solution}

\begin{solution}{exo:g7:light-sources-propagation:6}
\omterm{def:g7:light-sources-propagation:coats}{透明的}：窗玻璃、浅而清澈的水。\omterm{def:g7:light-sources-propagation:coats}{半透明的}：
描图纸、磨砂信封。\omterm{def:g7:light-sources-propagation:coats}{不透明的}：砖块、锡箔。
\end{solution}

\begin{solution}{exo:g7:light-sources-propagation:7}
亮度可以通过，而像被阻挡（打乱）。浴室的
窗户透进日光却不让外人看见；窗帘将
刺眼的点状灯丝变为宽广柔和的光晕。
\end{solution}

\begin{solution}{exo:g7:light-sources-propagation:8}
\omterm{def:g2:moon-in-the-sky:moon}{月球}上没有\omterm{def:g2:air-around-us:air}{空气}：宇航员头顶上没有任何东西把
经过的阳光向他们的眼睛漫射，因此天空
没有可见之物——正午也是黑的。地球的\omterm{def:g2:air-around-us:air}{空气}把
来自天空各方向的一部分阳光转发出来：我们明亮的蓝色穹顶
便是头顶上的漫射。
\end{solution}

\begin{solution}{exo:g7:light-sources-propagation:9}
每增加一层，光都会被打乱并多吸收一点：这些
类别取决于厚度。水是本章的第二个例子——
一杯水时透明，一片大海时不透明。（薄纸
是\omterm{def:g7:light-sources-propagation:coats}{半透明的}，合上的书是\omterm{def:g7:light-sources-propagation:coats}{不透明的}。）
\end{solution}

\begin{solution}{exo:g7:light-sources-propagation:10}
供\omterm{def:g6:light-propagation:ray}{光束}照射的东西：舞台上的薄雾或烟——
成群的细小\omterm{def:g7:light-sources-propagation:diffusion}{漫射物体}，接住\omterm{def:g6:light-propagation:ray}{光束}的光，把一部分
光散射到每个座位上。它们正在制造漫射，这是
侧面的眼睛看到\omterm{def:g6:light-propagation:ray}{光束}的唯一途径。
\end{solution}

\begin{solution}{exo:g7:light-sources-propagation:11}
火焰（\omterm{def:g7:light-sources-propagation:primary}{自身发光的光源}）$\to$ 沿直线到达脸部 $\to$ 脸
漫射 $\to$ 沿直线到达你的眼睛：链条中
只有一个\omterm{def:g7:light-sources-propagation:diffusion}{漫射物体}。经过湖面：脸漫射的光沿直线
到达静水，被镜面\omterm{def:g5:mirrors-reflection:reflection}{反射}（不是漫射——湖面是
抛光的），再沿直线到达你的眼睛：仍然只有一个
\emph{\omterm{def:g7:light-sources-propagation:diffusion}{漫射物体}}，外加一次镜面\omterm{def:g5:mirrors-reflection:reflection}{反射}——先是脸，再是
湖面那次规整的折返。
\end{solution}

\begin{solution}{exo:g7:light-sources-propagation:12}
阳光穿过\omterm{def:g2:air-around-us:air}{空气}时，沿路径一直有一部分被漫射
掉——路径越长，损失越大。\omterm{def:g1:day-and-night:daynight}{日落}时
太阳光沿最长的路径掠过\omterm{def:g5:air-pressure-weather:pressure}{大气层}，因此
直射\omterm{def:g6:light-propagation:ray}{光束}到达你眼中时已被夺走并减弱——那个温和的
橙色圆盘。被夺走的那部分并未消失：它
沿途被散射，成为其他天空的
日光——其中有一些在较晚时被散射，把高处的云染成
粉色。（哪种颜色最先被夺走，是下一年的内容。）
\end{solution}

\begin{solution}{pb:g7:light-sources-propagation:1}
\textbf{1.} 大理石本身不发光：看见需要光从
物体射入眼睛，不发光的\omterm{def:g7:light-sources-propagation:diffusion}{漫射物体}什么也送不出。
白色是对\emph{转发}能力的承诺，没有光送来
便毫无意义。
\textbf{2.} \omterm{def:g7:light-sources-propagation:primary}{自身发光的光源}：仅有灯具。\omterm{def:g7:light-sources-propagation:diffusion}{漫射物体}：
雕像、画作、参观者——还有展柜玻璃，微弱地
（大多数情况下它透射\omterm{def:g6:light-propagation:ray}{光线}）。
\textbf{3.} 漫射：大理石表面把接收到的光
向各个方向散射，因此大厅中每个位置都能得到光照，
且无人被刺眼。
\textbf{4.} 镜面式\omterm{def:g5:mirrors-reflection:reflection}{反射}：抛光黄铜以相等的角度把
灯光\omterm{def:g5:mirrors-reflection:reflection}{反射}到一个方向——刺眼的光
只落在那条单一路径所指之处。
\textbf{5.} \omterm{def:g7:light-sources-propagation:coats}{半透明的}：灯罩透过亮度
但打乱了像——灯丝的点状眩光变成
柔和的光球。
\textbf{6.} 依\omterm{def:g5:mirrors-reflection:reflection}{反射}规律，灯在清漆中显现，当
灯与观察者位于画作表面两侧，角度相等。
设计者把灯装得高而陡，使其\omterm{def:g5:mirrors-reflection:reflection}{反射}光路
俯冲到画前地面，
绝不射入主要观看位置的眼睛。
\textbf{7.} 玻璃大多透射\omterm{def:g6:light-propagation:ray}{光线}，但总会镜面\omterm{def:g5:mirrors-reflection:reflection}{反射}
微弱的一部分——\omterm{def:g7:light-sources-propagation:coats}{透明的}窗户是弱\omterm{def:g5:mirrors-reflection:reflection}{镜子}，在黑暗背景
的衬托下可以看见。黑色壁龛为透射景象提供
暗色底衬，使微弱的\omterm{def:g5:mirrors-reflection:reflection}{反射}无法与之竞争——幽灵般的
衣袖的鬼影随之淡去。
\textbf{8.} 两者皆是：从物理上说，这些丝线是尘埃微粒的漫射，
即\omterm{def:g6:light-propagation:ray}{光束}的光被散射到参观者的眼中；而尘土飞扬的
\omterm{def:g2:air-around-us:air}{空气}确实使其成为可能。物理归设计者，
诗意归策展人。
\textbf{9.} 没有\omterm{def:g1:light-and-shadows:source}{光源}，就没有射出的光，没有可接收的东西：
白色不过是对空无一物施加的高转发率。
条件一对所有物体同时不成立。
\textbf{10.} 大厅上方柔和、均匀、无影的日光——
只有亮度而无像：看不到云景，也没有
直射阳光在大理石上的刺眼光斑。
\textbf{11.} 天空本身：\omterm{def:g2:air-around-us:air}{空气}漫射阳光。北向光
是由半球形\omterm{def:g2:air-around-us:air}{空气}转发的阳光——柔和、
无\omterm{def:g1:light-and-shadows:source}{光源}，却仍属太阳；只有直射\omterm{def:g6:light-propagation:ray}{光束}
被挡在外。
\textbf{12.} 例如：“漫射是我们最好的朋友：它
把每件艺术品都送到每位参观者眼前，且毫无刺眼之感。
镜面式表面——清漆、黄铜、展柜玻璃——必须
留意，因为它们沿精确、
可预测的路径把灯光送入眼中。各类玻璃各司其职：\omterm{def:g7:light-sources-propagation:coats}{透明的}
用于透视，\omterm{def:g7:light-sources-propagation:coats}{半透明的}用于消减眩光和天光，\omterm{def:g7:light-sources-propagation:coats}{不透明的}用于
遮挡暗处。”
\end{solution}
